\documentclass[11pt]{article}
\usepackage[margin=1in]{geometry}
\usepackage{amsmath,amssymb,amsthm,booktabs,hyperref,microtype}
\hypersetup{colorlinks=true,linkcolor=blue,pdftitle={State-Weight Cocycles and Physical Length},pdfauthor={Yang--Mills representational dynamics research programme}}
\newtheorem{theorem}{Theorem}
\newtheorem{proposition}{Proposition}
\newtheorem{warning}{Warning}
\title{State-Weight Cocycles and Physical Length\\Target 57 of the H1 Downstream Programme}
\author{Yang--Mills representational dynamics research programme}
\date{21 July 2026}
\begin{document}
\maketitle
\begin{abstract}
We prove the scale-dependent state-weight cocycle for changing retained state
spaces.  Its contraction factors compose exactly and retain the initial root
moment.  We then distinguish contraction per descriptive RG level from decay
per physical distance.  A fixed contraction per geometrically growing level
is power-law in terminal scale, not a mass gap.  The correct bridge is a
length-weighted contraction cocycle.
\end{abstract}

\section{Changing-state-space cocycle}
Let $S_n$ be finite state spaces, $Q_n:S_n\to S_{n+1}$ positive subkernels,
and $V_n:S_n\to[1,\infty)$.
\begin{theorem}[Weight cocycle]
If
\begin{equation}\label{eq:step}
 Q_nV_{n+1}\leq q_nV_n,
\end{equation}
then
\[
 Q_0Q_1\cdots Q_{N-1}V_N
 \leq\left(\prod_{n=0}^{N-1}q_n\right)V_0.
\]
For a probability root law $\nu_0$, the total bad-chain mass is at most
$\nu_0(V_0)\prod_nq_n$.
\end{theorem}
\begin{proof}
Apply positivity successively to equation~\eqref{eq:step}.  Since $1\leq V_N$,
the terminal total mass is bounded by the weighted mass.
\end{proof}

The executable six-step star chain has exact mass
$2.0603224686684013\times10^{-16}$ and upper bound
$2.0603224686684008\times10^{-16}$, the difference being floating-point
roundoff.

\section{Physical-length cocycle}
Attach a positive physical increment $\ell_n$ to step $n$.  Define
\[
 m_N=\frac{-\sum_{n<N}\log q_n}{\sum_{n<N}\ell_n}.
\]
Uniform exponential rigidity in physical distance requires
\begin{equation}\label{eq:mass}
 \prod_{n<N}q_n\leq
 \exp\left(-m\sum_{n<N}\ell_n\right)
\end{equation}
with $m>0$ independent of cutoff, volume, and $N$.

\begin{proposition}[Geometric-level obstruction]
If $q_n=q<1$ and a terminal scale grows as $R_N=b^N$, then
\[
 q^N=R_N^{-\gamma},\qquad \gamma=-\frac{\log q}{\log b}.
\]
Thus a fixed contraction per RG level is naturally power-law in $R_N$.
\end{proposition}

For increments $\ell_n=2^n$ and the H1 value of $q$, the twenty-level effective
mass is only $1.1481773\times10^{-4}$ and tends to zero.  In contrast,
$q_n=e^{-m\ell_n}$ makes $m_N=m$ exactly; the executable witness gives
$m=0.4$ through total length $1023$.

\begin{warning}
The cocycle theorem transports representation-dependent rarity.  It becomes a
mass statement only after the increments are tied to physical separation and
after temporal observables are reconstructed densely.  Neither identification
is supplied by the finite kernel algebra.
\end{warning}

\section{Audit}
The rectangular cocycle is C2; equation~\eqref{eq:mass} is a C3 conditional
bridge.  Physical RG kernels, uniform length-scaled contraction, continuum OS
construction, and a Hamiltonian gap remain open.  Eleven hostile tests cover
changing dimensions, root moments, invalid kernels, and both scale laws.
\end{document}
